\clearpage
\fancyhead[L]{\small CONNAISSANCE DES TEMS.}
\SourceLeaf{620}
\begin{center}\small Göttingische gelehrte Anzeigen. Stück 191. Seiten 1905..1910. 29. November 1819.\\[3pt]\rule{.72\linewidth}{.3pt}\end{center}

\noindent\emph{Connaissance des tems, ou des mouvemens célestes, à l'usage des astronomes et des navigateurs, pour l'an {\upshape 1821}. Publiée par le bureau des longitudes.} Paris 1818. Bei der Witwe Courcier. 336 Seiten. Großoktav.

\begingroup\brokenpenalty=10000\relax
Der Anhang des vorliegenden Jahrgangs der Connaissance des tems zeichnet sich besonders durch mehrere gehaltreiche Aufsätze der führenden Geometer Frankreichs aus. Zuerst folgen die Ergebnisse einer theoretischen Untersuchung über die Libration des Mondes von \textsc{Poisson}. Hier werden auch neuere Beobachtungen des Mondflecks Manilius von \textsc{Bouvard} erwähnt, aus denen sich für die Neigung des Mondäquators $1^\circ\,27'\,40''$ ergebe (\textsc{Tobias Mayer} hatte 1749 aus Beobachtungen desselben Flecks $1^\circ\,29'$ gefunden). Mehr Beobachtungen dieser Art wären zweifellos sehr wünschenswert. --- Neue Methode zur Bestimmung der Kometenbahnen von \textsc{Lagrange}. Diese Abhandlung war bereits im Berliner astronomischen Jahrbuch für 1783 auf Deutsch erschienen. Die Methode setzt sechs vollständige Beobachtungen voraus, die paarweise nahe beieinanderliegen müssen. Ist die Lage der Bahnebene bekannt, so lässt sich aus dem geozentrischen Ort der heliozentrische ableiten, und zwei nahe beieinanderliegende heliozentrische Örter ergeben bekanntlich eine sehr einfache Bestimmung des halben Parameters. Daraus ist ersichtlich, dass jedes Paar nahe beieinanderliegender geozentrischer Örter eine Bedingungsgleichung zwischen der Bahnneigung, der Länge des aufsteigenden Knotens und dem halben Parameter ergibt. Die sechs vollständigen Beobachtungen ergeben daher drei solcher Bedingungsgleichungen, und durch Elimination von zwei der unbekannten Größen gelangt man schließlich zu einer Gleichung siebten Grades. Beim gegenwärtigen Stand dieses Teils der astronomischen Rechnung wird man keinen Anlass haben, diese Methode anzuwenden. --- Über die Rotation der Erde, von \textsc{Laplace}. Dieser Aufsatz enthält mehrere interessante Ausführungen über die freie Drehachse des Erdkörpers, sofern dieser mit einer Flüssigkeit bedeckt ist, die sich an ihrer Oberfläche im Gleichgewicht befindet, und über die Veränderung der Lage dieser Achse durch die Anziehung des Mondes und der Sonne. --- Den letzten Gegenstand behandelt die folgende Abhandlung von \textsc{Poisson} über die Präzession der Tagundnachtgleichen ebenfalls auf ähnliche Weise, mit der formalen Eleganz, die man von diesem ausgezeichneten Geometer gewohnt ist. --- Über den Einfluss der großen Gleichung von Jupiter und Saturn auf die Bewegung der Körper des Sonnensystems, von \textsc{Laplace}. Nur bei den eigenen Trabanten jener Planeten erweist sich dieser Einfluss als merklich. Dort zeigt er sich als eine Ungleichheit von ähnlicher Form wie die großen Gleichungen selbst. Beim vierten Jupiterstrabanten findet \textsc{Laplace} als größten Wert 14 Sekunden. --- Geografische Lage einer großen Zahl von Punkten im Mittelmeer, die der Fregattenkapitän \textsc{Gauttier} in den Jahren 1816 und 1817 bestimmt hat. --- Über das Gesetz der Schwere unter der Annahme, dass das Erdsphäroid homogen ist und dieselbe Dichte hat wie das Meer, von \textsc{Laplace}. Es ist bemerkenswert, dass unter dieser Annahme, die in der Natur freilich nicht zutrifft, die Schwere an der Erdoberfläche, sei es auf dem Meer, auf einer Insel oder auf dem Festland, von der Gestalt des über das Meer hinausragenden Festlands unabhängig wird. --- Nach einer späteren Bemerkung des Verfassers ist dieses Ergebnis unabhängig von der Annahme, dass das Erdsphäroid \emph{dieselbe} Dichte wie das Meer hat, und setzt \emph{nur} die gleichförmige Dichte des Erdsphäroids voraus. Und da die für die Schwere gefundene Formel vom Äquator zum Pol eine merklich geringere Zunahme ergibt als die Erscheinungen bei den Pendelschwingungen, ist damit bewiesen, dass die Dichte der Erde \SourceLeaf{621}von der Oberfläche zum Mittelpunkt hin zunimmt. Diese Zunahme der Dichte muss sich sogar bis in beträchtliche Tiefe erstrecken. --- Über die Positionen der Fixsterne, für die in den Ephemeriden die Monddistanzen angegeben werden, von \textsc{Bouvard}. Aus dem \textsc{Bessel}-\textsc{Bradley}-Katalog für 1755 und dem Katalog von \textsc{Piazzi} für 1800 (dessen Rektaszensionen jedoch, wie angegeben wird, anhand von Pariser Beobachtungen geringfügig geändert sind) werden die Positionen auf das Jahr 1830 übertragen. Diese Reduktion scheint jedoch nicht durchweg genau berechnet worden zu sein. --- Untersuchungen zu vier älteren Kometen, von \textsc{Burckhardt}. Für den zweiten Kometen von 1766 findet Herr \textsc{Burckhardt} unter Einbeziehung der Beobachtungen, die \textsc{Lanux} auf der Insel Bourbon gemacht hat, eine elliptische Bahn mit einer Umlaufzeit von 5 Jahren. Die Bewegungen der Kometen von 1774, 1723, 1729 lassen sich besser durch hyperbolische als durch parabolische Bahnen darstellen. --- Vom selben Astronomen eine Tabelle zur Bestimmung der Zeit bei parabolischer Bewegung aus zwei Abständen von der Sonne und dem eingeschlossenen Winkel. Außerdem kritische Bemerkungen zu verschiedenen Sternen, die \textsc{Bradley} nur einmal beobachtet hat. --- Über die englische Expedition zum Nordpol. --- Auszug aus der Beschreibung der Expedition von Kapitän \textsc{Tuckey} zur Erforschung des Flusses Zaire. --- Über die Gestalt der Erde und das Gesetz der Schwere an ihrer Oberfläche, von \textsc{Laplace}. Der hier abgedruckte Text gibt einen Überblick über die Hauptergebnisse einer Abhandlung, die der Verfasser in der Königlichen Akademie vorgetragen hat und die uns vollständig im Druck vorliegt. Darin wird die Aufgabe vollständiger und auf natürlichere Weise behandelt als je zuvor. Die Neuheit und Fruchtbarkeit der darin verwendeten Methoden und die Bedeutung der Ergebnisse für die Geologie machen sie zu einer der schönsten Arbeiten dieses großen Geometers. --- Auszug aus einem Brief von Herrn W. \textsc{Lambton} an Herrn \textsc{Delambre} über die ostindische Gradmessung, die damals (im September 1817) schon einen Meridianbogen von fast 10 Grad umfasste und noch immer weitergeführt wurde. Herr \textsc{Lambton} gibt die Breite des südlichen Endpunkts Runnä mit $8^\circ\,9'\,38''4$, die des nördlichen Endpunkts Darumgedda mit $18^\circ\,3'\,23''26$ und die Länge des Bogens mit 598610 Fathoms bei $62^\circ$ Fahrenheit an. Im Vergleich mit der französischen Gradmessung findet er die Abplattung $\frac{1}{309.5}$, mit der englischen $\frac{1}{313.54}$, mit der schwedischen $\frac{1}{307.15}$; aus dem Vergleich des ganzen Bogens mit dem ganzen Bogen von Dünkirchen bis Montjouy erhält er dagegen nach einer Formel von \textsc{Playfair} $\frac{1}{355}$ (diese starke Abweichung beruht auf einer fehlerhaften Umrechnung des französischen Bogens von Toisen in Fathoms). --- Auszug mit einigen Ortsbestimmungen in Schweden aus den Abhandlungen der Stockholmer Akademie. --- Über den Kometen von 1818, von \textsc{Nicollet}. --- Auszug aus \textsc{Lemonniers} Beobachtungsjournal, in dem sich nicht weniger als 12 Uranusbeobachtungen finden, die erste vom 14. Oktober 1750, die letzte vom 18. Dezember 1771. Im Jahr 1769 hatte dieser Astronom in 8 Tagen, vom 15. bis 23. Januar, den Planeten 6-mal beobachtet, ohne es zu wissen, und man begreift kaum, wie ihm selbst die Entdeckung entgangen sein kann. Herr \textsc{Bouvard}, der uns diese Beobachtungen bekannt macht, arbeitet an neuen Uranustafeln. Angesichts des Zustands des Instruments und der Uhr sowie der geringen Sorgfalt des Beobachters haben die Beobachtungen selbst allerdings keinen großen Wert. --- Über die Refraktion bei geringen Höhen, von \textsc{Groombridge}. Dies ist ein Nachtrag zu einer früheren Abhandlung dieses Astronomen, in der er für jeden einzelnen Stern nur die mittlere Abweichung von seinen Tafeln angegeben hatte. Hier werden auch die äußersten Unterschiede mitgeteilt. Das kann in gewissem Maße helfen, zu beurteilen, wie weit die Beobachtungen untereinander übereinstimmen. Doch für die wichtige Frage, über die die Meinungen der Astronomen bisher geteilt sind, ob man bei der Reduktion das äußere oder das innere Thermometer oder ein Mittel aus beiden verwenden sollte, wäre es wünschenswert, dass alle einzelnen Beobachtungen veröffentlicht würden. Außerdem sind die Beobachtungen von vier Solstitien von 1816---\SourceLeaf{622}1818 beigefügt. Das Mittel der beiden Bestimmungen der Schiefe der Ekliptik an den beiden Sonnensolstitien weicht hier nicht merklich vom Mittel aus den beiden Wintersolstitien ab, und das Mittel aus allen ergibt für den Zeitpunkt der Herbsttagundnachtgleiche von 1817 die mittlere Schiefe $23^\circ\,27'\,48''29$.
\par\clearpage\endgroup
